\section{研究设计}

% 本章通常 3000–5000 字，详细说明研究方案。
% 写清楚"用什么方法、怎么做、为什么这样做"。

\subsection{研究框架}

% 用文字或流程图描述整体研究框架。
% 如果有技术路线图，可以放在 figures/ 目录中引用：
% \begin{figure}[htbp]
%   \centering
%   \includegraphics[width=0.8\textwidth]{research-framework.png}
%   \bicaption{研究框架图}{Research Framework}
%   \label{fig:framework}
% \end{figure}

\subsection{数据来源}

% 说明数据的来源、采集方式、样本量和时间范围。
% 如果使用公开数据集，注明数据集名称和获取途径。
% 如果是问卷调查，说明问卷设计、发放和回收情况。

\subsection{研究方法}

% 详细介绍所用的分析方法或模型。
% 包括方法的基本原理、选择理由、参数设置等。
% 如果涉及数学公式，使用 equation 环境：
% \begin{equation}
%   y = \beta_0 + \beta_1 x_1 + \varepsilon
%   \label{eq:model}
% \end{equation}

\subsection{实施流程}

% 按时间或步骤说明研究的实施过程。
% 可以用 itemize 列出主要步骤。
